
%----------------------------------------------------------------------------------------
\chapter{The Technological Framework and Software Stack}
\label{chap:techstack}
\chaptermark{Technological Stack}
%----------------------------------------------------------------------------------------
\section[Backend Infrastructure]{Backend Infrastructure and Data Platform Architecture}

To operationalize the domain-oriented Data Product model and support interactive analytical workloads within Amadori S.p.A., 
the enterprise data architecture is undergoing an evolutionary transition from an on-premises relational infrastructure toward a scalable, 
cloud-native data platform. 

The baseline enterprise environment relies primarily on an on-premises Oracle database, which historically serves as the operational transaction 
repository and the target for analytical queries. To overcome the computational constraints of centralized on-premises storage and eliminate 
contention between transactional processing and resource-intensive analytical calculations, the target architecture introduces Google BigQuery 
as the enterprise cloud data warehouse. BigQuery provides a serverless, highly scalable environment based on distributed compute engines and 
columnar storage, enabling parallelized execution of complex analytical queries over large-scale historical datasets.

Rather than implementing an abrupt cutover, the transition follows a phased migration strategy. 
In the initial phase, analytical assets and ingestion pipelines continue to read directly from the operational Oracle database. 
As data models, historical partitions, and ETL routines are progressively refactored and validated, tables are transferred to Google BigQuery. 
During this transitional period, both database connections remain active simultaneously.

The presence of a centralized Semantic Layer implemented in Looker plays a strategic role in facilitating this dual-connection phase. 
Because Looker functions via in-database querying rather than data extraction, it can maintain concurrent dialect connections to both Oracle and BigQuery. 
Business dimensions and metrics formalized in LookML can initially target tables hosted on the Oracle instance; subsequently, 
as individual domains complete their migration to BigQuery, the underlying connection references within the LookML models are updated to point 
to the cloud tables. This decoupling ensures that downstream analytical consumers, executive reports, and AI interfaces remain completely shielded 
from the physical migration process, experiencing no disruption in metric definitions or query behavior while backend storage is shifted.

%----------------------------------------------------------------------------------------
% SECTION 5.2
%----------------------------------------------------------------------------------------

\section[The Semantic Layer with Looker and LookML]{Formalizing Business Data: The Semantic Layer with Looker and LookML}
\label{sec:looker_lookml}

To operationalize the conceptual definitions established by the enterprise ontology and provide an auditable single source of truth across analytical processes, 
the architecture implements Looker and its declarative modeling language, LookML. In this stack, Looker does not operate merely as a data visualization layer, 
but acts as the centralized execution engine for the enterprise Semantic Layer \cite{patel2023comparative}. 
Looker operates on an in-database query paradigm: analytical data remains stored within the underlying database, while the semantic platform dynamically compiles 
user interactions, dashboard requests, and API calls into optimized SQL queries executed directly by the target database engine at runtime \cite{gartner2021magic}.

This in-database approach plays a crucial role during the transition from the on-premises Oracle database toward Google BigQuery. 
Currently, the operational material planning pipeline populates the core fact table (\texttt{F\_MATERIAL\_PLANNING}) via a daily scheduled batch job running directly 
on the Oracle database. While the long-term target architecture aims to migrate these workloads to Google BigQuery, the abstraction provided by LookML allows the 
organization to manage this shift transparently. Because database connections are configured at the model level, LookML isolates physical connection dialects and 
schema differences. When tables are migrated to BigQuery in future deployment phases, connection strings and dialect syntax are updated within the model files without 
altering the exposed business entities, ensuring zero downtime and complete consistency for downstream users and cognitive agents.

\subsection{Architectural Best Practices and Performance Optimization}
Designing a high-throughput Semantic Layer requires balancing analytical flexibility with computational efficiency. 
When modeling complex supply chain dynamics in LookML, several engineering best practices were applied to minimize query execution overhead and prevent maintenance 
bottlenecks \cite{googlecloudlookerbp2024}:
\begin{itemize}
    \item \textbf{Pushing transformations to the backend:} Complex transformations, recursive allocations, and historical rolling aggregations are executed upstream 
    during database batch processing rather than calculated dynamically inside LookML. 
    Looker is dedicated to aggregation, dimensional filtering, and metric definition, preventing resource-intensive row-level computations at query runtime.
    \item \textbf{Denormalization of frequently filtered attributes:} While relational normalization minimizes redundancy in transactional platforms (OLTP), 
    analytical queries run faster when frequently joined descriptive attributes are pre-aggregated or denormalized into wide structures, reducing join depth.
    \item \textbf{Handling surrogate keys and natural codes:} Redundant internal surrogate keys are systematically hidden from the interface. 
    Instead, the model exposes standardized, concatenated business attributes (combining natural alphanumeric codes with their descriptions), 
    streamlining both human analysis and machine interpretation.
    \item \textbf{Minimizing runtime join overhead in dashboard visualizations:} While an Explore model file defines an extensive relational network to grant complete 
    self-service flexibility, individual dashboard tiles must query only the strictly necessary views. 
    Looker generates SQL dynamically based on the dimensions selected in each tile; keeping dashboard queries isolated to minimal join paths directly reduces runtime 
    join operations on the database engine, eliminating computational bottlenecks and query execution delays.
    \item \textbf{Curating the semantic exposure boundary:} Fields that are not explicitly declared in LookML remain completely invisible to consuming tools. 
    Restricting exposed dimensions to certified attributes establishes an architectural boundary that protects users from deprecated columns.
\end{itemize}

\subsection{Centralized Governance and the DRY Principle in Security Design}
A foundational engineering objective of formalizing business semantics in LookML is strict compliance with the \textit{Don't Repeat Yourself} (DRY) principle, 
which directly prevents code duplication and long-term maintenance overhead across enterprise analytics.

In conventional reporting architectures, enforcing data access segregation between distinct operational entities typically generates architectural anti-patterns. 
Technical teams are frequently forced to replicate entire models, SQL queries, or dashboard copies simply to hardcode facility-specific filtering criteria 
(e.g., maintaining cloned reports for each regional logistics deposit) or to strip out sensitive data columns for non-executive audiences. 
This duplication incurs severe operational risks: whenever an underlying business logic changes—such as updating the calculation formula of gross material requirements 
or shifting the baseline demand threshold—the alteration must be manually propagated across every replicated asset, creating high maintenance costs, security vulnerabilities, 
and eventual metric drift.

Within the LookML framework, operational business definitions and data security boundaries are decoupled, enabling the enterprise to enforce granular data access over a single,
unified codebase:
\begin{itemize}
    \item \textbf{Dynamic Row-Level Security (RLS) via \texttt{access\_filter}:} Rather than creating dedicated Explores or dashboards for different facilities, row-level 
    boundaries are enforced dynamically at query generation time. Through Looker's \texttt{access\_filter} parameter, the platform maps an authenticated user's session 
    attribute (e.g., \texttt{piattaforma\_autorizzata}) directly to the physical platform column (\texttt{l\_ptf.piattaforma\_cod}). Because Looker natively supports 
    multi-value user attributes and wildcards, this mechanism accommodates both local logistics operators restricted to a specific set of platforms and central planning 
    teams authorized to inspect the entire operational network. When a user runs a query or loads a dashboard, Looker dynamically injects the parameterized filtering 
    predicate into the SQL \texttt{WHERE} clause, ensuring strict data compartmentalization across the multi-site supply chain without duplicating a single line of 
    modeling code.
    \item \textbf{Granular Column-Level Security (CLS) and Access Grants:} Field visibility is governed systematically at the semantic declaration level rather than by 
    partitioning underlying database tables. Non-actionable technical fields—such as surrogate identifiers, system timestamps, and internal status flags—are universally 
    obscured from user visibility using \texttt{hidden: yes}, preventing clutter in the Explore interface and filtering unnecessary noise from the AI agent's semantic context. 
    Furthermore, to safeguard confidential commercial or financial metrics without generating separate reporting views, LookML provides \texttt{required\_access\_grants}. 
    For instance, while gross material volumes (\texttt{fabbisogno\_materiale}) are accessible to all operational planners, a related measure tracking standard component 
    unit costs (\texttt{costo\_standard\_materiale}) can be locked behind an access grant tied to an executive or managerial permission set. Users or agents lacking that 
    explicit entitlement are completely prevented from querying the restricted measure or seeing it in their semantic schema, ensuring security by design while maintaining 
    a single, centralized model.
\end{itemize}

\subsection{Implementation: Material Planning Fact and Master Views}
The operationalization of the Material Planning engine described in \Cref{chap:demandforecast} is implemented through a set of dedicated LookML views and explores. 
The primary fact view, \texttt{f\_material\_planning}, defines the daily gross material requirements across forecast dates, finished items, logistics platforms, and 
packaging or ingredient components:

\begin{lstlisting}[language=SQL, caption={LookML Definition for Material Planning Fact View}, label={lst:lookml_fact}]
view: f_material_planning {
  sql_table_name: `gsc-dp-mt-plan-sbx.material_planning_export.F_MATERIAL_PLANNING` ;;
  label: "Fabbisogno lordo materiali"

  dimension: primary_key {
    primary_key: yes
    type: string
    sql: concat(cast(${data_mp} as string), '-', cast(${art_id} as string), '-', 
                cast(${ptf_id} as string), '-', cast(${mat_id} as string), '-', 
                cast(${mat_um_id} as string)) ;;
    description: "Chiave primaria univoca: DATA_MP x ART_ID x PTF_ID x MAT_ID x MAT_UM_ID"
    hidden: yes
  }

  dimension_group: data_previsione {
    type: time
    label: "Data Previsione"
    timeframes: [raw, date, week, month, quarter, year]
    datatype: date
    convert_tz: no
    sql: parse_date('%Y%m%d', cast(${data_mp} as string)) ;;
    description: "Data di previsione del fabbisogno lordo"
    synonyms: ["data di previsione", "data fabbisogno", "orizzonte temporale", "data pianificata"]
  }

  measure: fabbisogno_materiale {
    type: sum
    label: "Fabbisogno Materiale"
    sql: ${mat_qta} ;;
    description: "Fabbisogno lordo del materiale (imballo, etichetta o ingrediente) necessario a coprire la produzione pianificata, per deposito, data e unita di misura"
    value_format_name: decimal_0
    synonyms: ["fabbisogno lordo", "quantita materiale", "consumo previsto", "fabbisogno totale"]
    drill_fields: [dettaglio_fabbisogno*]
  }

  set: dettaglio_fabbisogno {
    fields: [materiale.materiale, materiale.tipo_materiale, data_previsione_date, fabbisogno_materiale, um_cod]
  }
}
\end{lstlisting}

In this view, technical surrogate keys (\texttt{art\_id}, \texttt{mat\_id}, \texttt{ptf\_id}) and raw numeric dates are suppressed from user visibility using 
\texttt{hidden: yes}. Instead, the view exposes a clean \texttt{dimension\_group} parsing the integer date into standard calendar intervals. 
Business measures—such as \texttt{fabbisogno\_materiale}—incorporate pre-configured \texttt{drill\_fields}, enabling users to click on any aggregated volume to 
immediately inspect the underlying constituent articles and material classifications.

Master data views follow the same design. In \texttt{l\_art\_jde}, physical columns are enriched to dynamically distinguish between parent finished items, 
neutral items (\textit{PN}), and material components:

\begin{lstlisting}[language=SQL, caption={LookML Master Data View with Derived Categorization}, label={lst:lookml_master}]
view: l_art_jde {
  sql_table_name: `gsc-dp-mt-plan-sbx.material_planning_export.L_ART_JDE` ;;
  view_label: "Articolo"

  dimension: articolo {
    type: string
    label: "Articolo"
    sql: concat(${TABLE}.ART_JDE_LCOD, ' - ',${TABLE}.ART_JDE_DES) ;;
    description: "Codice e descrizione alfanumerica dell'articolo finito"
    synonyms: ["prodotto finito", "articolo di vendita", "prodotto", "PF", "PFL", "PFN"]
  }

  dimension: tipo_materiale {
    type: string
    label: "Tipo Materiale"
    sql: case 
           when ${TABLE}.ART_JDE_LCOD like 'I%' then 'Imballo'
           when ${TABLE}.ART_JDE_LCOD like '0%' then 'Ingrediente'
           else 'Altro'
         end ;;
    description: "Classificazione del materiale tra imballaggio e ingrediente primario"
    synonyms: ["categoria materiale", "tipologia componente", "imballo vs ingrediente"]
    suggestions: ["Imballo", "Ingrediente"]
  }
}
\end{lstlisting}

\subsection{Model Definition: Relational Joins and Access Perimeter}
The relational orchestration of the Material Planning domain is codified within the LookML model file (\texttt{material\_planning.model.lkml}). 
The model specifies the active database connection, imports the view definitions, and defines the join topology centered around the primary fact view:

\begin{lstlisting}[language=SQL, caption={LookML Model and Explore Configuration for Material Planning}, label={lst:lookml_model}]
connection: "materialplanningbq"

include: "/views/*.view.lkml"

explore: f_material_planning {
  label: "Fabbisogno lordo materiali"
  sql_always_where: 1=1 ;;

  # access_filter: {
  #   field: l_ptf.piattaforma_cod
  #   user_attribute: piattaforma_autorizzata
  # }

  join: articolo {
    from: l_art_jde
    view_label: "Articolo"
    type: inner
    relationship: many_to_one
    sql_on: ${f_material_planning.art_id} = ${articolo.art_jde_id} ;;
    fields: [articolo.articolo, articolo.articolo_attivo, articolo.articolo_neutro_cod]
  }

  join: materiale {
    from: l_art_jde
    view_label: "Materiale"
    type: inner
    relationship: many_to_one
    sql_on: ${f_material_planning.mat_id} = ${materiale.art_jde_id} ;;
    fields: [materiale.materiale, materiale.tipo_materiale]
  }

  join: l_ptf {
    type: inner
    relationship: many_to_one
    sql_on: ${f_material_planning.ptf_id} = ${l_ptf.ptf_id} ;;
  }

  join: l_um {
    type: inner
    relationship: many_to_one
    sql_on: ${f_material_planning.mat_um_id} = ${l_um.um_id} ;;
  }

  join: l_lin {
    type: inner
    relationship: many_to_one
    sql_on: ${articolo.art_jde_lin_id} = ${l_lin.lin_id} ;;
  }

  join: l_fam {
    type: inner
    relationship: many_to_one
    sql_on: ${articolo.art_jde_fam_id} = ${l_fam.fam_id} ;;
  }
}
\end{lstlisting}

As demonstrated in Listing~\ref{lst:lookml_model}, the Explore implements a role-playing dimension by joining the master table \texttt{l\_art\_jde} 
twice under two distinct semantic contexts: once as \texttt{articolo} (representing the finished commercial product generating demand) and once as 
\texttt{materiale} (representing the physical packaging or ingredient consumed). Furthermore, by explicitly filtering exposed attributes via the \texttt{fields} parameter, 
the model guarantees that only certified dimensions reach the analytical and conversational interfaces.

\subsection{Semantic Modeling for Self-Service Exploration and AI Enablement}
The enrichment of LookML parameters serves two interdependent consumption modes: human self-service analytics and cognitive AI agents.

For corporate planners, the centralized Explore defines certified relational joins across logistics platforms (\texttt{l\_ptf}), units of measure (\texttt{l\_um}), 
production lines (\texttt{l\_lin}), and the dual aliases of \texttt{l\_art\_jde} (mapped once as the parent \textit{Articolo} and once as the dependent \textit{Materiale}). 
This semantic foundation enables intuitive self-service BI: business users can build ad-hoc operational queries in the Explore interface, save customized views 
(\textit{Looks}), and combine them into interactive, cross-filtered dashboards without writing SQL code.

Simultaneously, this enriched metadata provides the deterministic guardrails required for autonomous Generative AI agents \cite{googlecloudlookerca2024}:
\begin{itemize}
    \item \textbf{Vocabulary Disambiguation via Synonyms:} Natural language queries rarely match exact database column names. The \texttt{synonyms} parameter maps 
    common enterprise jargon, informal acronyms (e.g., "PFL", "PFN", "deposito"), and domain terms directly to validated LookML dimensions, allowing an AI agent to match 
    user intent with deterministic schema fields.
    \item \textbf{Semantic Anchoring via Descriptions:} The \texttt{description} parameter defines the exact operational scope and mathematical logic of each metric. 
    This metadata prevents language models from misinterpreting raw values or fabricating fictitious aggregation logic.
    \item \textbf{Context Pruning via Hidden Fields:} By hiding raw IDs and technical join keys (\texttt{hidden: yes}), the Semantic Layer filters out noise from the 
    agent's context window. The LLM evaluates only certified, business-relevant fields, substantially reducing prompt token consumption and preventing hallucinated 
    joins \cite{googlecloudlookerca2024}.
\end{itemize}

By encapsulating business logic directly within version-controlled modeling files, Looker serves as an auditable, deterministic API. The operational orchestration of 
this semantic layer with autonomous conversational agents and large language models is formally analyzed and evaluated in \Cref{chap:conversational_agent}.

%----------------------------------------------------------------------------------------
% SECTION 5.3
%----------------------------------------------------------------------------------------

\section[Metadata Architecture Roadmap]{Governance Roadmap: Extending Lineage and Enterprise Knowledge via OpenMetadata}
\label{sec:openmetadata_roadmap}

While Looker operates as the centralized semantic compiler and enforces metric governance at the analytical layer, comprehensive visibility across the broader enterprise 
data lifecycle requires a dedicated metadata management platform. Currently, business definitions, dimensional attributes, and operational descriptions are codified and 
maintained directly within LookML files using parameters such as \texttt{description}, \texttt{label}, and \texttt{synonyms}. While this approach successfully establishes 
LookML as an executable "Business Glossary as Code", it remains primarily accessible to analytics engineers and technical practitioners working within the version-controlled 
modeling environment.

To decouple business glossary governance from technical codebases and achieve automated data lineage across heterogeneous systems, the enterprise architecture envisions 
the integration of OpenMetadata as an evolutionary roadmap milestone. OpenMetadata is an open-source, schema-first metadata management platform built upon standard JSON 
schemas for data discovery, governance, and quality profiling. Within the planned roadmap, OpenMetadata will expand enterprise governance beyond the analytical 
consumption boundary, focusing on three foundational capabilities:

\begin{enumerate}
    \item \textbf{Centralized, Low-Code Business Glossary:} Transitioning business definitions from plain-text LookML parameters into an intuitive enterprise catalog 
    accessible to non-technical business stakeholders \cite{dama2024dmbok}. This provides domain stewards across planning, supply chain, and commercial divisions with 
    the tooling required to manage term ownership, certify metric definitions, and approve terminology changes collaboratively without committing code to a Git repository.
    \item \textbf{Automated End-to-End Data Lineage:} Establishing systemic traceability across the entire data lifecycle. By connecting directly to source catalogs—tracing 
    tables from their operational origin in the on-premises Oracle database, through the cloud landing zones and transformations in Google BigQuery, down to the views, 
    dimensions, and measures compiled by Looker. Automated lineage tracking will enable data engineering teams to perform proactive impact analysis, accurately identifying 
    which downstream Explores, operational dashboards, or predictive material planning routines will be disrupted prior to executing any upstream schema migration.
    \item \textbf{Enterprise Knowledge Graph Construction:} Structuring enterprise metadata entities—including relational storage assets, batch pipeline schedules, 
    semantic metric definitions, and user access policies—as an interconnected knowledge graph \cite{hogan2021knowledge}. This structural context provides deterministic 
    insights into data provenance and reliability, serving as an authoritative grounding foundation to enhance future autonomous AI workflows.
\end{enumerate}

In this target framework, OpenMetadata will ingest metadata programmatically from Looker and underlying database instances via dedicated metadata APIs, eliminating 
manual synchronization overhead. This evolutionary enhancement strictly decouples operational query execution (handled in-database by Looker, Oracle, and BigQuery) 
from systemic metadata governance, completing the structural transition toward enterprise-grade Data Products.

\bigskip
In summary, the technological infrastructure deployed at Amadori S.p.A. harmonizes the operational stability of on-premises transactional databases with the analytical 
elasticity of cloud data warehousing. By establishing Looker as an in-database Semantic Layer, the enterprise enforces the DRY principle across its metric catalog, 
guarantees granular access governance, and encapsulates complex supply chain definitions within certified LookML models. With the storage infrastructure, relational modeling, 
and metadata governance roadmap formalized, the subsequent chapter explores how this unified semantic layer enables autonomous conversational analytics through Gemini in Looker.